\documentclass[11pt]{article}
\usepackage[T1]{fontenc}
\usepackage{lmodern}
\usepackage{amsmath,amssymb,amsthm}
\usepackage[margin=1.05in]{geometry}
\usepackage[hidelinks]{hyperref}

\newtheorem{theorem}{Theorem}
\newtheorem{lemma}{Lemma}
\newtheorem{corollary}{Corollary}
\newtheorem{remark}{Remark}

\newcommand{\C}{C}
\newcommand{\G}{G}
\newcommand{\N}{N}
\newcommand{\M}{M}

\title{Compact Groups Give an Arbitrary-Weight Partial Answer to the Dugundji Compact Question}
\author{Run \texttt{fa\_banach\_001}, agent \texttt{agent\_super\_01}}
\date{2026-07-01}

\begin{document}
\maketitle

\section*{Source Question}

Banakh and Kubis ask in arXiv:0610795 whether, for a Dugundji compact space
\(K\) of weight \(>\aleph_1\), the Banach space \(\C(K)\) must be
\(1\)-Plichko.  They prove the affirmative result at weight \(\leq\aleph_1\)
and note that compact groups are Dugundji compacta.

\section*{Result}

\begin{theorem}
Let \(G\) be a compact Hausdorff group.  Then \(\C(G)\) has a commutative
norm-one projectional skeleton.  In particular, \(\C(G)\) is \(1\)-Plichko.
\end{theorem}

\section*{Idea of the Proof}

Use the group structure to replace the ordinal inverse-sequence argument in
the source paper.  For every closed normal subgroup \(N\) such that \(G/N\)
is metrizable, Haar averaging over \(N\) is a norm-one projection of
\(\C(G)\) onto the functions factoring through \(G/N\).  These projections
commute because convolution of Haar probability measures on compact normal
subgroups \(N\) and \(M\) is Haar measure on \(NM\).  Peter--Weyl supplies
density, since finite-dimensional matrix coefficients factor through
metrizable quotients.

\section*{Proof}

Let \(\Gamma\) be the family of all closed normal subgroups \(N\triangleleft G\)
such that \(G/N\) is metrizable.  Order \(\Gamma\) by reverse inclusion:
\[
   N \preceq M \quad\Longleftrightarrow\quad N\supseteq M .
\]
This is directed because \(N\cap M\in\Gamma\), as \(G/(N\cap M)\) embeds as a
closed subgroup of the metrizable compact group \(G/N\times G/M\).

For \(N\in\Gamma\), let \(m_N\) denote normalized Haar measure on \(N\), and
define
\[
   (P_N f)(g)=\int_N f(gn)\,dm_N(n)\qquad (f\in \C(G),\ g\in G).
\]
Then \(P_N\) is a contractive projection.  It has norm one unless \(G\) is
empty, and its range is exactly the subspace of functions constant on right
cosets of \(N\), i.e. the canonical copy of \(\C(G/N)\).  Thus every range is
separable.

If \(N\supseteq M\), then every \(N\)-coset-constant function is
\(M\)-coset-constant, and direct averaging gives
\[
   P_NP_M=P_N=P_MP_N .
\]
For arbitrary \(N,M\in\Gamma\), the convolution \(m_N*m_M\) is the normalized
Haar measure on the compact subgroup \(NM\).  Indeed it is supported on
\(NM\), and because \(N\) and \(M\) are normal it is left invariant under both
\(N\) and \(M\), hence under the subgroup generated by them, namely \(NM\).
Therefore
\[
   P_NP_M=P_{NM}=P_MP_N ,
\]
so the projections commute.

Now let \(N_1\succeq N_2\succeq\cdots\) be an increasing sequence in
\(\Gamma\), equivalently \(N_1\supseteq N_2\supseteq\cdots\), and set
\(N=\bigcap_j N_j\).  Then \(G/N\) embeds into \(\prod_j G/N_j\), hence is
metrizable, so \(N\in\Gamma\).  The quotient maps \(G/N\to G/N_j\) separate
points.  By the complex Stone--Weierstrass theorem, the union of the pulled
back algebras \(\bigcup_j \C(G/N_j)\) is dense in \(\C(G/N)\).  In terms of
the projections, this says
\[
   P_N[\C(G)] = \overline{\bigcup_j P_{N_j}[\C(G)]}.
\]
This is the required countable-supremum axiom for a projectional skeleton.

It remains to check that the union of all ranges is dense in \(\C(G)\).  By
the Peter--Weyl theorem, finite linear combinations of matrix coefficients of
finite-dimensional continuous unitary representations of \(G\) are uniformly
dense in \(\C(G)\).  If \(\pi:G\to U(n)\) is such a representation, then every
matrix coefficient factors through \(G/\ker\pi\), and \(G/\ker\pi\) is
metrizable because it is isomorphic to the compact subgroup \(\pi(G)\subseteq
U(n)\).  Hence every matrix coefficient lies in \(P_{\ker\pi}[\C(G)]\) for
some \(\ker\pi\in\Gamma\).  Thus \(\bigcup_{N\in\Gamma}P_N[\C(G)]\) is dense.

The family \((P_N)_{N\in\Gamma}\) is therefore a commutative norm-one
projectional skeleton on \(\C(G)\).  By the standard characterization of
\(1\)-Plichko spaces as Banach spaces with a commutative norm-one
projectional skeleton, \(\C(G)\) is \(1\)-Plichko.

\section*{Scope and Novelty Check}

The result covers compact groups of arbitrary weight, hence an important
Dugundji compact subclass, but it does not settle the source question for all
Dugundji compacta.

Local run indexes and parsed sources were searched for compact-group
Plichko variants.  The local hits record the source paper's weight
\(\aleph_1\) compact-group application and later references saying that
\(\C(K)\) is \(1\)-Plichko for Abelian compact groups.  I found no existing
run packet or local source hit stating the above all-compact-group
projectional-skeleton argument.

\section*{References}

\begin{itemize}
\item T. Banakh and W. Kubis, \emph{Spaces of continuous functions over
Dugundji compacta}, arXiv:0610795.
\item W. Kubis, \emph{Banach spaces with projectional skeletons}, J. Math.
Anal. Appl. 350 (2009), 758--776.  Used for the projectional-skeleton
characterization of Plichko spaces as quoted in local later sources.
\item The Peter--Weyl theorem for compact groups.
\end{itemize}

\section*{Human Review Recommendation}

Check the nonabelian averaging step \(P_NP_M=P_{NM}\), specifically the
claim that \(m_N*m_M\) is Haar measure on \(NM\) for compact normal
subgroups.  This is the only group-theoretic point where commutativity of the
group is not available.

\end{document}
